% Chapter 20: Building Better Sentences: From Simple to Skilful

\chapter{Building Better Sentences: From Simple to Skilful}

\section{Pedagogical Purpose}
Learners now enter Unit 7, which shifts focus from individual words to whole sentences. Having built strong vocabulary and grammar accuracy, they are ready to learn how to make their sentences richer, more varied, and more precise — key skills for effective writing.

\begin{learningobjectives}
The learner can:
\begin{itemize}
    \item expand a short sentence with additional relevant detail.
    \item combine two short sentences into one longer sentence using a connecting word.
    \item identify and remove repeated or unnecessary words in a sentence.
    \item arrange words into the correct, natural sentence order.
    \item replace a vague word with a more precise one.
    \item write a short paragraph with varied sentence length.
\end{itemize}
\end{learningobjectives}

\section{Prerequisite Check}
\begin{quickcheck}
Is this a complete sentence? "The dog." (\textit{No — it has no verb; it is not a complete thought.})
\end{quickcheck}

\begin{rememberbox}
In Chapter 1, you learned that a sentence needs a subject and a verb to express a complete thought. Now you will learn how to make your sentences more interesting and skilful.
\end{rememberbox}

\section{Chapter Opener}
Rohan has written a set of very short sentences for his class diary, and Ms. Sharma reads them aloud.

\textbf{Rohan's Diary:} "I went to the park. I saw a dog. The dog was big. The dog barked. I was scared. I ran home."

\textbf{Ms. Sharma:} "Rohan, every sentence here is correct, but reading six short sentences in a row feels choppy, like riding over small bumps! Let's learn how to smooth this out and build stronger sentences."

\textbf{Maya:} "Can we join some of them together?"

\textbf{Ms. Sharma:} "Exactly, Maya! Let's explore several ways to build better sentences."

\section{Language Experience}
\textbf{Two Versions of Rohan's Diary}

\begin{tcolorbox}[colback=gray!5, colframe=gray!50, sharp corners, boxsep=0pt, left=5pt, right=5pt, top=5pt, bottom=5pt, breakable]
\textbf{Choppy:} I went to the park. I saw a dog. The dog was big. The dog barked. I was scared. I ran home.
\end{tcolorbox}

\begin{tcolorbox}[colback=gray!5, colframe=gray!50, sharp corners, boxsep=0pt, left=5pt, right=5pt, top=5pt, bottom=5pt, breakable]
\textbf{Improved:} I went to the park, where I saw a huge dog. When it barked loudly, I was so scared that I ran all the way home.
\end{tcolorbox}

\section{Look and Notice}
\begin{thinkbox}
\begin{enumerate}
    \item How many separate sentences are in the choppy version? How many in the improved version?
    \item What extra details were added to the improved version (like "huge" and "loudly")?
    \item Which words were used to join the short sentences together?
    \item Which version is more interesting to read, and why?
\end{enumerate}
\end{thinkbox}

\section{Discover / Explicit Teaching}
Skilful writers use several tools to build better sentences:

\begin{table}[h!]
\centering
\begin{tabular}{|l|l|l|}
\hline
\textbf{Tool} & \textbf{What it does} & \textbf{Example} \\
\hline
Expansion & Adds detail (when, where, how, why) & "The dog barked." $\rightarrow$ "The big dog barked loudly at the gate." \\
Combination & Joins short sentences with connecting words & "I was tired. I went to bed." $\rightarrow$ "I was tired, so I went to bed." \\
Removing repetition & Cuts unnecessary repeated words & "The red red apple" $\rightarrow$ "The red apple" \\
Word order & Arranges words naturally & "Quickly ran the boy" $\rightarrow$ "The boy ran quickly" \\
Precision & Replaces vague words with exact ones & "The thing was nice" $\rightarrow$ "The gift was thoughtful" \\
\hline
\end{tabular}
\end{table}

\section{Grammar Word}
\begin{grammarword}{SENTENCE EXPANSION}
\textbf{Sentence expansion} means adding extra, relevant detail to a simple sentence to make it more informative and interesting, usually by answering questions like \textit{when, where, how,} or \textit{why}.

\textit{Examples:}
\begin{itemize}
    \item Simple: "The bird sang."
    \item Expanded: "Early in the morning, the small bird sang sweetly from the tree."
\end{itemize}

\textit{Non-example:} Adding unrelated or unnecessary detail is not true expansion: "The bird, which reminds me of my lunch, sang." (irrelevant detail)

\textit{Quick Check:} Expand this sentence with one detail of "when": "The children played." \textit{(e.g., "After school, the children played.")}
\end{grammarword}

\section{Core Explanation}

\subsection*{Combining short sentences}
``It was raining. We stayed inside.'' $\rightarrow$ ``It was raining, \textbf{so} we stayed inside.'' \textit{Check:} Does the connecting word correctly show how the two ideas relate (reason, contrast, time)?

\subsection*{Removing repetition}
``The tall, tall boy climbed the tall tree.'' $\rightarrow$ ``The tall boy climbed the tree.'' \textit{Check:} Is the repeated word adding meaning, or just cluttering the sentence?

\subsection*{Correct word order}
``Quickly the girl the ball threw.'' (incorrect order) $\rightarrow$ ``The girl quickly threw the ball.'' (natural order). \textit{Check:} Does the sentence sound natural when read aloud?

\subsection*{Precision over vagueness}
``The thing was nice.'' $\rightarrow$ ``The gift was thoughtful.'' / ``The weather was nice.'' $\rightarrow$ ``The weather was warm and sunny.'' \textit{Check:} Could I picture this clearly if I only heard this word?

\section{Worked Example}
\begin{examplebox}
\textbf{Sentences to improve:} "The cat sat. The cat was on the mat. The mat was soft."

\textit{Step 1:} Identify the repeated subject --- "the cat" and "the mat" appear more than once.
\textit{Step 2:} Combine the ideas into one sentence, keeping only necessary detail.
\textit{Step 3:} Choose a natural word order and remove repetition.

\textbf{Answer:} "The cat sat on the soft mat."
\end{examplebox}

\section{Guided Practice}
\begin{activitybox}{Activity A: Identification}
Find the repeated or unnecessary word in each sentence.
\begin{enumerate}
    \item "The big, big elephant walked slowly."
    \item "The good good news made everyone happy."
\end{enumerate}
\end{activitybox}

\begin{activitybox}{Activity B: Completion}
Combine each pair of sentences using a suitable joining word.
\begin{enumerate}
    \item "It was cold. I wore a jacket."
    \item "Maya likes painting. Rohan likes cricket."
\end{enumerate}
\end{activitybox}

\section{Independent Practice}
\begin{activitybox}{Independent Practice}
\begin{enumerate}
    \item Expand this sentence with two details (when and where): "The children laughed."
    \item Combine these three short sentences into one: "The sun was hot. We were thirsty. We drank water."
    \item Replace the vague word "nice" with a more precise word in: "We had a nice time at the fair."
    \item \textit{Reasoning:} Why do you think using only very short sentences, one after another, can make writing feel dull, even if every sentence is grammatically correct?
\end{enumerate}
\end{activitybox}

\section{Grammar Detective}
\begin{detectivebox}
\textbf{Student sentence:}

``The thing was good. It was good because it was a good thing to have.''

\textit{Questions:}
\begin{enumerate}
    \item What is wrong?
    \item What should it be?
    \item Why?
\end{enumerate}

\textit{Guidance:} The sentence is vague and repetitive ("thing," "good" repeated). Improve with specific words: "The new bicycle was wonderful because it let me ride to school faster."
\end{detectivebox}

\section{Common Confusion}
\begin{warningbox}
Learners sometimes think that combining sentences means simply joining them with "and" every time, which can create a long, unclear "and... and... and" chain instead of a well-structured sentence.

\textit{How to check:} Ask, "Does 'and' really show the right relationship, or would 'but,' 'so,' 'because,' or 'when' fit better?"
\end{warningbox}

\section{Writing Connection}
Rewrite this choppy paragraph as 3-4 improved sentences, using expansion, combination, and precise words: "I went to the shop. I bought a pen. The pen was blue. I was happy. I came home."

\section{Summary}
\begin{summarybox}
\begin{itemize}
    \item \textbf{Expansion} adds relevant detail (when, where, how, why) to a simple sentence.
    \item \textbf{Combination} joins short sentences using suitable joining words.
    \item \textbf{Removing repetition} keeps sentences clean and clear.
    \item \textbf{Correct word order} makes sentences sound natural.
    \item \textbf{Precision} replaces vague words with exact ones.
    \item Skilful writers \textbf{vary} their sentence length and structure.
\end{itemize}
\end{summarybox}

\section{Key Terms}
\begin{table}[h!]
\centering
\begin{tabular}{|l|l|}
\hline
\textbf{Term} & \textbf{Simple Meaning} \\
\hline
Sentence expansion & Adding relevant detail to a simple sentence \\
Sentence combination & Joining short sentences into one using a joining word \\
Repetition & Unnecessary repeating of a word or idea \\
Word order & The natural arrangement of words in a sentence \\
Precision & Choosing exact, specific words \\
Sentence variety & Using different sentence lengths and structures \\
\hline
\end{tabular}
\end{table}

\begin{selfcheck}
I can:
\begin{itemize}
    \item expand a simple sentence with relevant detail.
    \item combine short sentences using a suitable joining word.
    \item remove unnecessary repetition from a sentence.
    \item arrange words in correct, natural order.
    \item replace vague words with precise ones.
    \item vary sentence length in my own writing.
\end{itemize}
\end{selfcheck}